% =============================================================================
% BEAMER SLIDE TEMPLATE
%
% HOW TO USE:
%   Replace every \ph{NAME} call with real content.
%   \ph{} renders as bold [NAME] in the compiled PDF — easy to spot.
%   Also replace the color values on lines marked COLOR below.
%   Uncomment \includegraphics lines once the figure files exist.
%
%   Figures go in a figures/ subdirectory next to this file.
% =============================================================================

\documentclass[t,pdf,aspectratio=169]{beamer}

\usepackage[activate={true,nocompatibility}]{microtype}
\usepackage{lmodern}
\usepackage{amsmath}
\usepackage{amsfonts}
\usepackage{breqn}
\usepackage{nicefrac}
\usepackage{dirtytalk}
\usepackage{caption}
\usepackage{listings}
\usepackage{tikz}
\usepackage{bm}
\usepackage{bbm}
\usepackage{amsmath,amssymb,amsopn,dsfont,mathrsfs,bm}

% ---- Placeholder helper (delete once all placeholders are filled) -----------
\newcommand{\ph}[1]{[#1]}

% ---- Theorem environments ---------------------------------------------------
\newtheorem{thm}{Theorem}
\newtheorem{pf}{Proof}
\newtheorem{dfn}{Definition}

\BeforeBeginEnvironment{definition}{%
    \setbeamercolor{block title}{fg=white,bg=PRIMARY}
    \setbeamercolor{block body}{fg=black, bg=PRIMARY!20!white}
}
\AfterEndEnvironment{definition}{
    \setbeamercolor{block title}{use=structure,fg=white,bg=structure.fg!75!black}
    \setbeamercolor{block body}{parent=normal text,use=block title,bg=block title.bg!10!bg}
}
\BeforeBeginEnvironment{theorem}{%
    \setbeamercolor{block title}{fg=white,bg=PRIMARY}
    \setbeamercolor{block body}{fg=black, bg=PRIMARY!20!white}
}
\AfterEndEnvironment{theorem}{
    \setbeamercolor{block title}{use=structure,fg=white,bg=structure.fg!75!black}
    \setbeamercolor{block body}{parent=normal text,use=block title,bg=block title.bg!10!bg}
}
\AtBeginEnvironment{thm}{%
    \setbeamercolor{block body}{bg=white, fg=black}
    \setbeamercolor{block title}{bg=PRIMARY, fg=white}
}

% ---- Appendix frame counter -------------------------------------------------
\newcommand{\backupbegin}{
   \newcounter{finalframe}
   \setcounter{finalframe}{\value{framenumber}}
}
\newcommand{\backupend}{
   \setcounter{framenumber}{\value{finalframe}}
}

% ---- Math helpers -----------------------------------------------------------
\newcommand{\indep}{\perp\!\!\!\perp}
\newcommand{\ind}[1]{\mathds{1}\left\{#1\right\}}

% ---- Theme and colors -------------------------------------------------------
\usetheme{Madrid}
\usecolortheme{default}

% COLOR: primary brand color as RGB in [0,1] — e.g. (64,92,126) -> (0.251,0.361,0.494)
\definecolor{PRIMARY}{rgb}{0.25098039, 0.36078431, 0.49411765}
% COLOR: lighter accent — e.g. (165,192,231) -> (0.647,0.753,0.906)
\definecolor{PRIMARY-light}{rgb}{0.647058823529412, 0.752941176470588, 0.905882352941176}

\setbeamercolor*{palette primary}{bg=PRIMARY, fg=white}
\setbeamercolor*{palette secondary}{bg=white, fg=PRIMARY}
\setbeamercolor*{palette tertiary}{bg=PRIMARY-light, fg=black}
\setbeamercolor*{palette quaternary}{bg=PRIMARY, fg=white}
\setbeamercolor{section number projected}{bg=PRIMARY}
\setbeamercolor{button}{fg=white, bg=PRIMARY-light}

\hypersetup{colorlinks, citecolor=PRIMARY, filecolor=PRIMARY,
            linkcolor=PRIMARY, urlcolor=PRIMARY}

\mode<presentation>{}
\setbeamertemplate{itemize item}{\color{black}$\bullet$}
\setbeamertemplate{itemize subitem}{\color{black}\scriptsize$\bullet$}

% =============================================================================
% TITLE BLOCK
% =============================================================================

\title[\ph{Short title}]{\ph{Full title}}
\subtitle{\ph{Subtitle}}          % remove line if no subtitle

\author[\ph{Short author(s)}]{
  \textbf{\ph{Author One}}\inst{1} \and
  \textbf{\ph{Author Two}}\inst{1} \and
  \textbf{\ph{Author Three}}\inst{2}
  % add/remove authors as needed
}
\institute[]{
  \inst{1} \ph{Institution One} \\
  \inst{2} \ph{Institution Two}
  % add/remove institutions as needed
}

\date[\ph{Short date}]{\ph{Venue} \\ \ph{Full date, e.g. April 8, 2026}}

% =============================================================================
\begin{document}
% =============================================================================

\begin{frame}
  \titlepage
\end{frame}

% Uncomment for table of contents:
% \begin{frame}{Overview}
%   \tableofcontents
% \end{frame}

% =============================================================================
\section{Introduction}
% =============================================================================

% --- Section title card ------------------------------------------------------
\begin{frame}
  \vfill
  \begin{center}
    \color{PRIMARY}{\huge{Introduction}}
  \end{center}
  \vfill
\end{frame}

% --- Research question -------------------------------------------------------
\begin{frame}{Introduction}
  \framesubtitle{Research Question}
  \vfill
  \begin{center}
    \textcolor{PRIMARY}{
      \Large
      \textbf{\ph{Research question line 1} \\
      \medskip \ph{Research question line 2}}
    }
  \end{center}
  \vfill
\end{frame}

% --- Motivation --------------------------------------------------------------
\begin{frame}{Introduction}
  \framesubtitle{Motivation}
  \vfill
  \begin{itemize}
    \item \ph{Motivation item 1}
    \item \ph{Motivation item 2}
    \item \ph{Motivation item 3}
    \item \ph{Motivation item 4}
  \end{itemize}
  \vfill
\end{frame}

% --- Optional: motivating figure (two-panel) ---------------------------------
% \begin{frame}{Introduction}
%   \framesubtitle{\ph{Figure slide subtitle}}
%   \vfill
%   \begin{center}
%   \begin{minipage}{0.49\textwidth}
%     \begin{figure}[htb]
%       \centering
%       \includegraphics[width=7cm]{figures/FIGURE_LEFT}
%     \end{figure}
%   \end{minipage}
%   \begin{minipage}{0.49\textwidth}
%     \begin{figure}[htb]
%       \centering
%       \includegraphics[width=7cm]{figures/FIGURE_RIGHT}
%     \end{figure}
%   \end{minipage}
%   \end{center}
%   \vfill
% \end{frame}

% --- Contribution ------------------------------------------------------------
\begin{frame}{Introduction}
  \framesubtitle{Contribution}
  \vfill
  We want to:
  \begin{itemize}
    \item \ph{Contribution 1}
    \item \ph{Contribution 2}
    \item \ph{Contribution 3}
    \item \ph{Contribution 4}
    \item \ph{Contribution 5}
  \end{itemize}
  \vfill
\end{frame}

% --- Literature --------------------------------------------------------------
\begin{frame}{Introduction}
  \framesubtitle{Literature}
  \vfill

  \textcolor{PRIMARY}{\textbf{\ph{Literature topic 1}}}
  \scriptsize
  \begin{itemize}
    \item[-] \ph{Author(s) (Year, Journal) -- description}
    \item[-] \ph{Author(s) (Year, Journal) -- description}
    \item[-] \ph{Author(s) (Year, Journal) -- description}
  \end{itemize}
  \medskip

  \normalsize
  \textcolor{PRIMARY}{\textbf{\ph{Literature topic 2}}}
  \scriptsize
  \begin{itemize}
    \item[-] \ph{Author(s) (Year, Journal) -- description}
    \item[-] \ph{Author(s) (Year, Journal) -- description}
  \end{itemize}
  \medskip

  \normalsize
  \textcolor{PRIMARY}{\textbf{\ph{Literature topic 3}}}
  \scriptsize
  \begin{itemize}
    \item[-] \ph{Author(s) (Year, Journal) -- description}
    \item[-] \ph{Author(s) (Year, Journal) -- description}
  \end{itemize}

  \vfill
\end{frame}


% =============================================================================
\section{Methodology}
% =============================================================================

\begin{frame}
  \vfill
  \begin{center}
    \color{PRIMARY}{\huge{Methodology}}
  \end{center}
  \vfill
\end{frame}

% --- Pipeline overview figure ------------------------------------------------
\begin{frame}{Methodology}
  \framesubtitle{\ph{Method overview subtitle}}
  \hypertarget{method-back}{}
  \begin{figure}[htb]
    \centering
    % \includegraphics[width=10cm]{figures/METHOD_OVERVIEW_FIGURE}
  \end{figure}
  \vfill
  % \hyperlink{method-detail}{\beamergotobutton{Details}}
\end{frame}

% --- Core math / notation ----------------------------------------------------
\begin{frame}{Methodology}
  \framesubtitle{\ph{Method detail subtitle}}
  \vfill

  \ph{Introductory text for the key equation.}

  \begin{equation}
    \text{\ph{KEY\_EQUATION}}   % replace with actual LaTeX math
  \end{equation}

  \ph{Text after the equation, e.g. what each term means.}

  \vfill
\end{frame}

% --- Challenges (all highlighted) --------------------------------------------
\begin{frame}{Methodology}
  \framesubtitle{Challenges}
  \vfill

  \textbf{\textcolor{PRIMARY}{\ph{Challenge 1 title}}}
  \scriptsize
  \begin{itemize}
    \item[-] \ph{Challenge 1, item 1}
    \item[-] \ph{Challenge 1, item 2}
    \item[-] \ph{Challenge 1, item 3}
  \end{itemize}
  \medskip
  \normalsize

  \textbf{\textcolor{PRIMARY}{\ph{Challenge 2 title}}}
  \scriptsize
  \begin{itemize}
    \item[-] \ph{Challenge 2, item 1}
    \item[-] \ph{Challenge 2, item 2}
  \end{itemize}
  \medskip
  \normalsize

  \textbf{\textcolor{PRIMARY}{\ph{Challenge 3 title}}}
  \scriptsize
  \begin{itemize}
    \item[-] \ph{Challenge 3, item 1}
    \item[-] \ph{Challenge 3, item 2}
  \end{itemize}
  \vfill
\end{frame}

% --- Theorem -----------------------------------------------------------------
\begin{frame}{Methodology}
  \framesubtitle{\ph{Theorem slide subtitle}}
  \vfill
  \begin{theorem}[\ph{Theorem name}]
    \ph{Theorem statement}
  \end{theorem}
  \vfill
  \ph{Discussion / construction that follows the theorem.}
  \vfill
\end{frame}

% --- Illustrative figure -----------------------------------------------------
\begin{frame}{Methodology}
  \framesubtitle{\ph{Figure subtitle}}
  \vfill
  \begin{figure}[htb]
    \centering
    % \includegraphics[width=10.5cm]{figures/METHOD_FIGURE}
    % \caption*{{\scriptsize Source: \ph{Source citation}}}
  \end{figure}
  \vfill
\end{frame}

% --- Simulation / power curve ------------------------------------------------
\begin{frame}{Methodology}
  \framesubtitle{\ph{Power curve subtitle}}
  \vfill
  \begin{figure}[h]
    \centering
    % \includegraphics[width=0.5\linewidth]{figures/POWER_FIGURE}
    % \caption{\ph{Caption for the power curve.}}
  \end{figure}
  \textcolor{PRIMARY-light}{\textit{Note: \ph{Details on simulation setup.}}}
  \vfill
\end{frame}

% --- Alternative approaches --------------------------------------------------
\begin{frame}{Methodology}
  \framesubtitle{Other Approaches}
  \vfill
  \textcolor{PRIMARY}{\textbf{\ph{Approach 1 title}}}
  \scriptsize
  \begin{itemize}
    \item[-] \ph{Approach 1, item 1}
    \item[-] \ph{Approach 1, item 2}
    \item[-] \ph{Approach 1, item 3}
  \end{itemize}
  \bigskip
  \normalsize
  \textcolor{PRIMARY}{\textbf{\ph{Approach 2 title}}}
  \scriptsize
  \begin{itemize}
    \item[-] \ph{Approach 2, item 1}
    \item[-] \ph{Approach 2, item 2}
  \end{itemize}
  \vfill
\end{frame}


% =============================================================================
\section{Results}
% =============================================================================

\begin{frame}
  \vfill
  \begin{center}
    \color{PRIMARY}{\huge{Results}}
  \end{center}
  \vfill
\end{frame}

% --- Main results figure -----------------------------------------------------
\begin{frame}{Results}
  \framesubtitle{\ph{Results subtitle}}
  \vfill
  \begin{figure}[h]
    \centering
    % \includegraphics[width=0.7\linewidth]{figures/RESULTS_FIGURE}
    % \caption{\ph{Results caption.}}
  \end{figure}
  \vfill
\end{frame}

% --- Two-column comparison ---------------------------------------------------
\begin{frame}{Results}
  \framesubtitle{\ph{Comparison subtitle}}
  \vfill
  \begin{columns}[t]
    \begin{column}{0.48\textwidth}
      \textbf{\ph{Group / condition 1}}
      \begin{itemize}
        \item[-] \ph{Item 1}
        \item[-] \ph{Item 2}
        \item[-] \ph{Item 3}
      \end{itemize}
    \end{column}
    \begin{column}{0.48\textwidth}
      \textbf{\ph{Group / condition 2}}
      \begin{itemize}
        \item[-] \ph{Item 1}
        \item[-] \ph{Item 2}
        \item[-] \ph{Item 3}
      \end{itemize}
    \end{column}
  \end{columns}
  \vfill
\end{frame}


% =============================================================================
\section{Next Steps}
% =============================================================================

\begin{frame}
  \vfill
  \begin{center}
    \color{PRIMARY}{\huge{Next Steps}}
  \end{center}
  \vfill
\end{frame}

\begin{frame}{Next Steps}
  \framesubtitle{\ph{Next steps subtitle}}
  \vfill
  \textcolor{PRIMARY}{\textbf{\ph{Next step 1 title}}}
  \scriptsize
  \begin{itemize}
    \item[-] \ph{Item 1}
    \item[-] \ph{Item 2}
    \item[-] \ph{Item 3}
  \end{itemize}
  \medskip
  \normalsize
  \textcolor{PRIMARY}{\textbf{\ph{Next step 2 title}}}
  \scriptsize
  \begin{itemize}
    \item[-] \ph{Item 1}
    \item[-] \ph{Item 2}
  \end{itemize}
  \medskip
  \normalsize
  \textcolor{PRIMARY}{\textbf{\ph{Next step 3 title}}}
  \scriptsize
  \begin{itemize}
    \item[-] \ph{Item 1}
    \item[-] \ph{Item 2}
  \end{itemize}
  \vfill
\end{frame}

% --- Closing slide -----------------------------------------------------------
\begin{frame}
  \vfill
  \begin{center}
    \color{PRIMARY}{\huge{Thank you for your attention!}}
  \end{center}
  \vfill
\end{frame}


% =============================================================================
% APPENDIX — frame numbers reset; do not count past the main deck
% =============================================================================

\backupbegin

\begin{frame}
  \vfill
  \begin{center}
    \color{PRIMARY}{\huge{Appendix}}
  \end{center}
  \vfill
\end{frame}

% --- Appendix: concept / background ------------------------------------------
\begin{frame}{Appendix}
  \framesubtitle{\ph{Appendix slide 1 subtitle}}
  \hypertarget{appendix-1}{}
  \vfill
  \begin{minipage}{0.49\textwidth}
    \begin{itemize}
      \item \ph{Item 1}
      \item \ph{Item 2}
      \item \ph{Item 3}
    \end{itemize}
  \end{minipage}
  \begin{minipage}{0.49\textwidth}
    \begin{figure}[htb]
      \centering
      % \includegraphics[width=7cm]{figures/APPENDIX_1_FIGURE}
      % \caption*{{\scriptsize Source: \ph{Source}}}
    \end{figure}
  \end{minipage}
  \vfill
  \hyperlink{appendix-1-back}{\beamergotobutton{Back}}
\end{frame}

% --- Appendix: math / theory 1 -----------------------------------------------
\begin{frame}{Appendix}
  \framesubtitle{\ph{Appendix slide 2 subtitle}}
  \hypertarget{appendix-2}{}
  \vfill

  \textcolor{PRIMARY}{\textbf{\ph{Theorem / problem name:}}}
  \ph{Introductory text.}

  \begin{equation}
    \text{\ph{EQUATION}}   % replace with actual LaTeX math
  \end{equation}

  \ph{Follow-up text or additional equations.}

  \vfill
  \hyperlink{appendix-2-back}{\beamergotobutton{Back}}
\end{frame}

% --- Appendix: definition ----------------------------------------------------
\begin{frame}{Appendix}
  \framesubtitle{\ph{Appendix slide 3 subtitle}}
  \hypertarget{appendix-3}{}
  \vfill
  \begin{definition}[\ph{Definition name}]
    \ph{Definition body.}
  \end{definition}
  \vfill
  \hyperlink{appendix-3-back}{\beamergotobutton{Back}}
\end{frame}

% --- Appendix: example texts (two-column) ------------------------------------
\begin{frame}{\ph{Example slide title}}
  \begin{columns}[t]
    \begin{column}{0.48\textwidth}
      \textbf{\ph{Group 1 label}}
      \begin{itemize}
        \item \ph{Example text from group 1.}
      \end{itemize}
    \end{column}
    \begin{column}{0.48\textwidth}
      \textbf{\ph{Group 2 label}}
      \begin{itemize}
        \item \ph{Example text from group 2.}
      \end{itemize}
    \end{column}
  \end{columns}
\end{frame}

\backupend

\end{document}
